\cotefile{12}
\scanpage{18}
\mspage{16}
Considérons $\Gamma=\underline{\mathrm{Aut}}_{\mathbb Q}(V_0)$, groupe algébrique % L1
sur $\mathbb Q$, et les deux hom % L2
\[
\deleted{\unsure{j_1,j_2}}i_1,i_2:\ G_{m\,\mathbb C}\longrightarrow\Gamma_{\mathbb C} % L3
\]
définissant la bigraduation de $V$ % L4
\[
\begin{aligned}&\deleted{\unsure{g_{\lambda}(x)|i_1(\lambda)}}\\[-2pt]&i(\lambda,\mu)\unsure{\longrightarrow|\mapsto|\text{est}}\lambda^p\mu^q\quad\text{sur }V^{pq}\end{aligned} % L5
\]
Soit $G$ le plus petit ss-groupe algébrique de \inserted{$\Gamma$} % L6
\deleted{tel} que $G_{\mathbb C}$ contienne $\mathrm{Im}\,i_1,\mathrm{Im}\,i_2$ ; % L7
\deleted{Alors} c'est aussi le ss-groupe algébrique de % L8
\deleted{\unsure{$G$}} $\Gamma_{\mathbb C}$ \deleted{\unsure{$G$}} engendré par les $i_1^\sigma$, $\sigma$ parcourant % L9
les \deleted{\unsure{en}} autom de $\mathbb C$. C'est un groupe % L10
algébrique connexe, appelé \uline{groupe de Mumford-Tate} % L11
de la structure de Hodge. \deleted{\unsure{Pour toute} \unsure{à quelque} % L12
représentation de $G$ \unsure{qui soit|qui} \unsure{par} sur un vectoriel % L13
$M$ de rang fini sur $\mathbb Q$, on} Si $\Sigma$ est % L14
une opération tensorielle en une variable, % L15
de sorte que $G$ opère sur $\Sigma(V_0)$, % L16
on voit que les élémts de $\Sigma(V_0)$ invariants % L17
sous $G$ sont \deleted{\unsure{exactt}} ceux qui sont % L18
invariants par \inserted{\unsure{les $i$.|les ss.|les homs}} $i_1,i_2$, i.e. qui sont de % L19
bidegré $0,0$. Cela permet de voir que % L20

% archive sheet number: 17
\scanpage{19}
\mspage{17}
le ``groupe de Galois'', au sens structures de % L1
Hodge, de $(V,V_0,(V^{pq}))$, s'identifie à $G$. % L2

Supposons maintenant que $V$ provienne d'% L3
un motif $M$, engendrant une \inserted{$\otimes$-}catégorie de % L4
motifs $C_0$. On a alors $C_0\unsure{-|\longrightarrow}C$, donc % L5
$G\longrightarrow G_0$, où $G_0$ est le groupe de Galois % L6
motivique ($G_0\subset\Gamma$). La conjecture de % L7
Hodge s'interprète en disant que \uline{$G=G_0$ !} % L8

14) \uline{Structure de positivité} % L9

% archive sheet number: 18
\scanpage{21}
\hfill\fbox{Motifs} % L1

1) La catégorie $M^+(X)$ % L2

2) Variance \unsure{avec|pour} $X$ % L3

3) Cas $X=\varprojlim X_i$ % L4

4) Foncteurs $T_\ell$ % L5

5) Les $\mathbb Q_\ell(-n)$ % L6

6) La catégorie $M(X)$ % L7

7) Les foncteurs $\mathrm{Hom}$ et $\mathrm{RHom}$ % L8

8) Motifs constants, tordus, et \unsure{pl. con.|pt. con.|pot. con.}. % L9

9) Filtrations de $M(X)$ et $M^+(X)$ % L10

10) Anneaux $M^+(X)$ et $M(X)$. % L11

11) Interprétation topologique des \deleted{\unsure{types de niveau}} \inserted{niveaux} % L12
par la filtration canonique \deleted{\unsure{de leurs}} variétés proj. non singulièr\unsure{es|e}. % L13

12) L'hom $L(K)\to M^+(K)$ et invariants birationnels fondamentaux. % L14
[N.B. À généraliser à des coefficients quelconques] % L15

13) Caractérisation galoisienne des filtrations (conj. de Tate) % L16

14) Invariants de Galois et \unsure{thm.|thms} de commutation. % L17

15) Cohomologie absolue des motifs. Cas des schémas de type fini. % L18

16) Relation avec les pts rationnels et la coh. des VA sur % L19
un schéma de type fini. % L20

17) Formes positives. % L21

18) Motifs et fonctions $L$ % L22

19) Relations avec la cohomologie de Hodge ....  % L23

% archive sheet number: 20
\scanpage{23}
\uline{Motifs} % L1

1) À tout préschéma \textsuperscript{noeth} (éventuellement de type fini sur un anneau % L2
noethérien) $X$ est associé une \uline{catégorie abélienne} % L3
$\mathcal M^+(X)$, dite catégorie des \uline{motifs} \inserted{\uline{effectifs}} sur $X$. C'est une % L4
$\mathbb Q$-catégorie \deleted{de $\mathbb Q$-vectoriels} abélienne, i.e. pour tt % L5
$M\in\mathrm{Ob}\,\mathcal M^+(X)$, et $n\in\mathbb Z$, $n\ne0$, $n\,\mathrm{id}_M$ est un isom % L6
de $M$. De plus, $\mathcal M^+(X)$ est muni d'un % L7
produit tensoriel commutatif et unitaire, \inserted{exact à droite,} % L8
% slanted
\margin{On peut en termes du \unsure{dernier|premier|données} construire des foncteurs $\unsure{\lambda^i|\bigwedge^i}M$ etc...} % L8a
l'unité est notée $\mathbb 1_X$ ou $\mathbb Q_X(0)$. % L9

On considère aussi la catégorie dérivée bornée % L10
$D^b(\mathcal M^+(X))$ de $\mathcal M^+(X)$. Le produit tensoriel est % L11
étendu en un bifoncteur $M\mathbin{\underline{\underline{\otimes}}}N$ en $M,N\in\mathrm{Ob}\,D^b(\mathcal M^+(X))$ % L12

2) Soit $f:X\longrightarrow Y$ un morphisme de préschémas % L13
\deleted{\unsure{loc}} noeth, il lui est associé un \uline{foncteur} % L14
\uline{exact} % L15
\[
f^*:\mathcal M^+(Y)\longrightarrow\mathcal M^+(X) % L16
\]
compatible avec $\otimes$, \inserted{d'où $\mathbb L f^*:D^{b}(\mathcal M^+(Y))\unsure{-|\to}D^{b}(\mathcal M^+(X))$} On a transitivité. % L17

Si $f$ est de type fini, et propre \uline{ou} $Y$ excellent, % L18
\unsure{On|on|on a} \deleted{\unsure{définit|défini}} de $=$ un foncteur % L19
\[
\mathbb R f_*:D^b(\mathcal M^+(X))\unsure{-|\longrightarrow}D^b(\mathcal M^+(Y)) % L20
\]

% archive sheet number: 22
\scanpage{24}
satisfaisant une formule de transitivité, % L1
et formule de projection % L2
% slanted
\margin{Considérer aussi la formule de dualité entre $\mathbb R f_*,f^*$, et les foncteurs $f^!,\mathbb R f_!$ et leurs relations d'adjonction et enfin les formules de \unsure{Künneth|chgt de base|standard} reliant tous ces foncteurs .....} % L2a
\[
\mathbb R f_*(M\otimes\mathbb L f^*(N))\simeq\mathbb R \unsure{f|f_*}(M)\otimes N % L3
\]

3) Supposons $X=\varprojlim X_i$, système projectif % L4
filtrant \unsure{en. affine|ess. affine}. Alors par les foncteurs % L5
images inverses, on a % L6
\[
\mathcal M^+(X)\simeq\underset{i}{\underset{\longrightarrow}{\mathrm{Lim}}}\,\mathcal M^+(X_i) % L7
\]

En particulier, si $X$ est de type fini % L8
sur $S=\mathrm{Spec}(A)$, alors $X$ est limite de % L9
préschémas $X_i$ de type fini sur $\mathbb Z$, et % L10
la détermination de $\mathcal M^+(X)$ avec ses % L11
structures déjà envisagées est \deleted{\unsure{dét|dit}} \unsure{ramené|ramenée} % L12
au cas des préschémas de type fini. % L13

De $=$, si $(S_i)$ est système % L14
projectif filtrant \unsure{en. affine|ess. affine}, $S=\varprojlim S_i$. % L15
et si $X,Y$ de t. f. sur $S$ sont % L16
définis par $(X_i)$, $(Y_i)$ de la % L17
façon habituelle, si on part de % L18

% archive sheet number: 23
\scanpage{25}
$M_{i_0}\in\mathrm{Ob}\,D^b(\mathcal M^+(X_{i_0}))$, d'où $M_i,M$, on % L1
aura pour $f_{i_0}:X_{i_0}-Y_{i_0}$ la relation % L2
\[
\mathbb R f_*(M)=\varinjlim_i v_i^*(\mathbb R f_{i*}(M_i)) % L3
\]
où $v_i:Y-Y_i$ est le morphisme % L4
canonique. % L5

4) Soit $\ell\in\mathbb P$ tel que $\ell\,1_X\in\Gamma(X,\mathscr O_X)$ soit % L6
inversible. Alors on a un foncteur % L7
\[
T_\ell=T_\ell^{(X)}:\mathcal M^+(X)\longrightarrow\mathcal M_\ell^+(X) % L8
\]
où $\mathcal M_\ell^{*}(X)$ est \inserted{la catégorie} \unsure{formé|formée} des ``$\mathbb Q_\ell$-modules % L9
constructibles sur $X$'', i.e. la catégorie % L10
déduite des la catégorie des ``systèmes % L11
$\ell$-adiques de faisceaux de $\ell$-torsion % L12
constructibles'' en négligeant précisément % L13
les faisceaux de torsion. Le foncteur % L14
$T_\ell$ est \inserted{\uuline{compatible avec} $\otimes$ \uuline{et unité}} \uuline{exact} et \fbox{\uuline{fidèle}} (mais non % L15
pleinement fidèle), \uuline{compatible avec le} % L16

% archive sheet number: 24
\scanpage{26}
\uuline{changement de base} $f^*$, \uuline{et compatible} % L1
\uuline{également avec} $\mathbb R f_*$ [N.B. $T_\ell$ s'étend % L2
% slanted
\margin{compatibilité avec $f^!$, $\mathbb R f_!$, avec l'hom résidu etc...} % L2a
évidemment en un foncteur % L3
$D^b(\mathcal M(X))\unsure{-|\to}D^b(\mathcal M_\ell(X))$]. $\{$ La détermination % L4
des $T_\ell$ est encore ramenée au cas où % L5
$X$ est de t.f. sur $\mathbb Z$. % L6

% slanted
\margin{\uline{NB.} Ceci exclut déjà le choix trivial $\mathcal M^+(X)=0$ pour tt $X$, car il faudrait qu'on ait $\mathbb R^{\unsure{*|x}}f_*(\mathbb Z_{\underline{\ell}})=0$ pour $f,\ell$, ce qui n'est pas vrai en général ---} % L7
\deleted{Si $X$ est un $\mathbb Q$-schéma, \inserted{de \deleted{\unsure{présentation}} \inserted{type} fini,} alors on % L7a
a aussi un foncteur % L8
\[
T_\infty=T_\infty^{(X)}:D^b(\mathcal M(X))\longrightarrow\mathcal M_\infty(X) % L9
\]
$\mathcal M_\infty(X)$ est la catégorie des \inserted{$\mathscr O_X$-}Modules % L10
cohérents. Ce foncteur $T_\infty$ est encore % L11
\uline{exact} et \uline{fidèle}, compatible avec % L12
$f^*$ et $\mathrm Rf_*$.} % L13
\uline{Signaler aussi la compatibilité de} $T_\ell$ \uline{avec l'isom. de Künneth} % L13a

5) Pour tout $X$, on a un \deleted{\unsure{un}} % L14
élément canonique % L15
\[
\deleted{\mathbb Z}\mathbb Q_X(-1)\ \text{ou}\ \mathbb 1_X(-1)\ \in\mathrm{Ob}\,\mathcal M^+(X) % L16
\]

% archive sheet number: 25
\scanpage{27}
dont la formation est compatible avec chgt % L1
de base (il suffit donc de le considérer % L2
sur $\mathrm{Spec}\,\mathbb Z$), avec \deleted{des} isom % L3
\[
T_\ell(\deleted{\mathbb Z}\mathbb Q(-1))\simeq\deleted{\mathbb Z}\mathbb Q_\ell(-1)={\underbrace{T_\ell(G_m)}_{\substack{\text{Module de Tate}\\\text{\inserted{inverse} de }G_m}}}^{-1} % L4
\]
et le cas échéant ($X$ sur $\mathbb Q$). % L5

On peut définir $\deleted{\mathbb Z}\mathbb Q(-1)$ comme % L6
$\mathrm R^2 f_*(\deleted{\mathbb Z}\mathbb 1_{P^1_X})$, où $f:P_X^1\longrightarrow X$ est la % L7
projection canonique. Posant % L8
\[
\deleted{\mathbb Z}\mathbb Q(-n)=\deleted{\mathbb Z}\mathbb Q(-1)^{\otimes n}\qquad\text{pour }n\geq0 % L9
\]
on peut prouver, à l'aide des % L10
axiomes déjà posés, que si $f:X-S$ % L11
est \inserted{lisse} projectif à fibres \inserted{géom.} connexes non vides, % L12
\deleted{partout} de dim relative $d$, alors % L13
\[
\mathbb R^{2d}f_*(\mathbb 1_X)\simeq\mathbb Z_S(-d) % L14
\]
et si on enlève l'hyp ``projectif'' mais % L15
seult $f$ quasi-projectif, on trouve encore % L16
\[
\mathbb R^{2d} f_!(\unsure{1|\mathbb 1}_X)\simeq\mathbb Z_S(-d) % L17
\]

% archive sheet number: 26
\scanpage{28}
On \deleted{exige} \inserted{veut} de plus \deleted{compatibilité}, \deleted{\unsure{avec|que}} si $X/S$ est % L1
lisse et $Y\underset{i}{\subset}X$ est lisse sur $S$, de codim % L2
$d$ dans $X$, l'isom % L3
\[
\mathbb R i^!(1_X)\simeq\deleted{\mathbb Z}\mathbb Q_Y(-d) % L4
\]
compatible avec les isom déjà connus % L5
du point de vue $\ell$-adique .... % L6

6) \deleted{Chaque $\mathcal M(X)$ est plongé dans} % L7
Le foncteur % L8
\[
M\unsure{\leadsto|\mapsto}M(-1)=M\otimes\deleted{\mathbb Z}\mathbb Q(\unsure{+1|-1}) % L9
\]
\deleted{\unsure{est|de}} $\mathcal M^+(\unsure{X|A})$ dans lui-$=$ est \uline{pleinement} \uline{fidèle} % L10
mais pas une équivalence en général. % L11
Il y a donc une façon canonique % L12
d'élargir $\mathcal M^+(X)$ en \deleted{\unsure{une}} $\mathcal M(X)$ de telle façon % L13
que \deleted{\unsure{de}} $\otimes\deleted{\mathbb Z}\mathbb Q(-1)$ devienne une équivalence, % L14
\deleted{\unsure{ce}}, en prenant la \inserted{pseudo-}limite inductive % L15
% slanted
\margin{$\mathcal M^+(X)$ est une \uline{sous-catégorie abélienne épaisse de} $\mathcal M(X)$ ; [\unsure{l'appart.} à $\mathcal M^+(X)$ se vérifie fibre par fibre ----]} % L15a
des % L16
\diagram{
\begin{tikzcd}[ampersand replacement=\&]
\mathcal M^+ \arrow[r,"{\otimes\deleted{\mathbb Z}\mathbb Q(-1)}"] \& \mathcal M^+ \arrow[r,hook,"{\otimes\deleted{\mathbb Z}\mathbb Q(-1)}"] \& \mathcal M^+ % L16a
\end{tikzcd}
}
. On prolonge % L16b
$\mathcal M(X)$ la structure $\otimes$, alors $\deleted{\mathbb Z}\mathbb Q(-1)$ % L17
devient inversible, soit $\deleted{\mathbb Z}\mathbb Q(1)$ son inverse, et % L18
tt élément de $\mathcal M(X)$ peut s'écrire ss forme % L19

% archive sheet number: 27
\scanpage{29}
$n$ assez grand, \unsure{sous} la forme % L1
$M_n(n)$, avec $M_n\in\mathrm{Ob}\,\mathcal M^+(X)$, [pour $n$ fixé, % L2
$M_n$ est bien déterminé par $M$ à isom unique % L3
près, et $\bigl(M_n(n)\simeq M_{n+1}(n+1)\bigr)\Leftrightarrow\bigl(M_{n+1}\simeq M_n(-1)\bigr)$, % L4
on retrouve la description de $\mathcal M(X)$ en % L5
termes de Pseudo $\underset{\longrightarrow}{\mathrm{Lim}}$. ]. Les foncteurs % L6
$T_\ell$ s'étendent à $\mathcal M(X)$, de façon \unsure{ess.|can.} unique, de façon % L7
à rester compatible avec $\otimes$. % L8

7) Dans $\mathcal M(X)$, on a aussi des foncteurs % L9
$\underline{\mathrm{Hom}}$, liés à $\otimes$ par la formule % L10
habituelle % L11
\[
\begin{aligned}
\mathrm{Hom}(P\otimes Q,R)&\simeq\mathrm{Hom}(P,\underline{\mathrm{Hom}}(Q,R))\\ % L12
&\simeq\mathrm{Hom}(Q,\underline{\mathrm{Hom}}(P,R)), % L13
\end{aligned}
\]
et \deleted{\unsure{satisfaisant|satisfaisant par suite}} \inserted{qui s'étendent en} % L14
\[
\mathbb R\mathrm{Hom}(\deleted{\unsure{Q}}P,Q)\qquad P,Q\in\mathrm{Ob}\,D^bM(X) % L15
\]
satisfaisant à la relation analogue % L16
relativement à $\underline{\otimes}$. \uuline{La formation des} % L17
$\underline{\mathrm{Hom}}$ \inserted{et $\mathbb R\underline{\mathrm{Hom}}$} \uuline{est compatible avec les} $T_\ell$. % L18

\deleted{NB.} [NB. On retrouve la formation % L19
des $f^!$ .... ] % L20

% archive sheet number: 28
\scanpage{30}
8) Soient $\ell,\ell'\in\mathbb P$ premiers aux car rés. % L1
de $X$. Soit $M\in\mathrm{Ob}\,\mathcal M^+(X)$. On veut % L2
que % L3
\[
\begin{gathered}
T_\ell(M)\ \text{faisceau constant tordu}\\ % L4
\Updownarrow\\ % L5
T_{\ell'}(M)\ \text{faisceau constant tordu} % L6
\end{gathered}
\]
et que sous ces conditions, $T_\ell(M)$ et % L7
$T_{\ell'}(M)$ doivent avoir $\unsure{\bar m|=}$ rang en chaque % L8
pt. / Pour vérifier l'égalité des rangs, on % L9
est ramené au cas où $X$ est \inserted{le} spectre d'un % L10
corps (fini si on veut, \deleted{\unsure{d}} ou clôture alg. % L11
d'un tel). Plus généralement, si % L12
$u$ est un endomorphisme de $M$, \deleted{\unsure{on|et}} % L13
($X=\mathrm{Spec}\,k$, $k$ un corps), on en déduit % L14
\[
T_\ell(u)\in\mathrm{End}\,T_\ell(M),\qquad T_{\ell'}(u)\in\mathrm{End}\,(T_{\ell'}(M)) % L15
\]
et je dis que on a % L16
\[
\boxed{\mathrm{Tr}\,T_\ell(u)=\mathrm{Tr}\,T_\ell(\unsure{u'|u})\in\deleted{\unsure{\mathbb Z}}\mathbb Q} % L17
\]
d'où, remplaçant $u$ par $\bigwedge^i u$, le fait % L18
\[
\boxed{P(T_\ell(u),t)=P(T_{\ell'}(u),t)\in\mathbb Q[t]} % L19
\]

% archive sheet number: 29
\scanpage{31}
[où il s'agit des polynômes caractéristiques] % L1

Pour ceci, notons que % L2
\[
u\in\mathrm{Hom}(\mathbb 1_X,\mathrm{Hom}(M,\unsure{M?|M_2|M}))=\mathrm{Hom}(\mathbb 1_X,\overset{\vee}{M}\otimes M) % L3
\]
et on a un morphisme \uline{contraction} % L4
\[
\overset{\vee}{M}\otimes M\longrightarrow\mathbb 1_X % L5
\]
% slanted
\margin{\uline{à inclure dans le formalisme tensoriel}} % L5a
d'où un $c(u)\in\mathrm{Hom}(\mathbb 1_X,\mathbb 1_X)$ % L6
et \deleted{\unsure{donc}} $\mathrm{Tr}\,T_\ell(u)=T_\ell(c(u))\in\deleted{\mathbb Z}\mathbb Q_\ell$. % L7
et il suffit de savoir : % L8
\[
\boxed{X\ \text{connexe}\Longrightarrow\mathrm{Hom}(\mathbb 1_X,\unsure{\mathbb 1|1}_X)\simeq\mathbb Q\,\mathrm{id}_{\mathbb 1_X}} % L9
\]

9) \uline{Filtration de} $\mathcal M^+(X)$ \uline{et} $\mathcal M(X)$ % L10

1°) \uline{Filtration par le poids}, \deleted{\unsure{et}} % L11
\[
\deleted{\unsure{\mathcal M^+(X)}}\ \mathcal M^{+0}(X)\subset\mathcal M^{+1}(X)\subset\cdots\subset\mathcal M^{+i}(X)\subset\cdots % L12
\]
\uline{filtration exhaustive} de $\mathcal M^+(X)$. [L'appartenance % L13
à $\mathcal M^{+i}(X)$ se vérifie fibre par fibre, \deleted{\unsure{l'app.}} % L14
% slanted
\margin{indices en bas} % L14a
\inserted{\deleted{\unsure{ou}} (géométrique si on veut)} et dans le cas $X$ de t.f. \inserted{sur $\mathbb Z$} il suffit de % L15
vérifier en les pts fermés. (Pour \unsure{voir|avoir|ceux-} % L16
\unsure{ça|ce|ci}, il y a un critère par Frobenius, cf % L17

% archive sheet number: 30
\scanpage{32}
plus bas).] \deleted{loc. pour} % L1
\deleted{$\mathcal M^i(X)$ est ``engendré'' par les % L1a
$\mathrm R^i f_{\underline *}$. L'appartenance à $\mathcal M^i(\unsure{X|k(x)})$ \unsure{---} de % L2
$M|_x$ implique l'appart\unsure{en.|enance}} Compatible avec % L3
$f^*$, et \deleted{l'appart} $M|_x\in\mathrm{Ob}\,\mathcal M^{+i}(k(x))$ % L4
$\Longrightarrow\exists\,U\ni x$ voisinage de $\bar x$, tel que % L5
$M|_U\in\mathrm{Ob}\,\mathcal M^{+i}(U)$. Soit $f:X-Y$, % L6
% slanted
\margin{\unsure{généraliser|quasi-limité|quasi-lisse} Compatibilité avec $\varprojlim$ des préschémas ....} % L6a
alors $\mathbb R^j f_!:\mathcal M^{+i}(X)\longrightarrow\mathcal M^{i+j}(Y)$ ; de % L7
plus, $\mathcal M^{+i}(k)$ ($k$ un corps alg. clos) % L8
est ``engendré'' par les \inserted{sous-espaces} $\mathrm R^j f_!(\deleted{\mathbb Z}\mathbb 1_X)$ pour % L9
$X$ si on veut projectif lisse de dim $\leq j$ sur $k$ % L10
[au sens que tout $M\in\mathrm{Ob}\,\mathcal M^i(k)$ a une % L11
filtration dont les facteurs sont isomorphes % L12
à de tels sous-espaces ]. % L13
% slanted
\margin{indices en bas} % L13a
\[
\deleted{\mathcal M^i(X)\otimes\mathbb Z(-1)} % L14
\]
\[
\left\{\begin{aligned}
&\mathcal M^{+i}(X)\otimes\mathcal M^{+j}(X)\subset\mathcal M^{+i+j}(X)\\ % L15
&\mathbb Z(-1)\in\mathrm{Ob}\,\mathcal M^{+\unsure{2|1}}(X) % L16
\end{aligned}\right.
\]
d'où % L17
\[
\mathcal M^{+i}(X)\otimes\mathbb Z(-j)\subset\mathcal M^{+i+2j}(X)\qquad\text{pour }j\geq0 % L18
\]

% archive sheet number: 31
\scanpage{33}
on a mieux : % L1
\[
M\in\mathrm{Ob}\,\mathcal M^{+i}(X)\Longleftrightarrow M(-j)\in\mathrm{Ob}\,\mathcal M^{+i+2j}(X) % L2
\]

De cette façon, la filtration \inserted{de $\mathcal M^+(X)$} par les % L3
$\mathcal M^{+i}(X)$ \inserted{$(i\geq0)$} peut se prolonger en une filtration % L4
de $\mathcal M(X)$ par des $\mathcal M^i(X)$ $(i\in\mathbb Z)$ de % L5
telle façon que pour $i\geq0$, on \unsure{suit|sent|soit} la filtration % L6
déjà \deleted{\unsure{$\mathcal M^+(X)$|$\mathcal M(X)$}} envisagée, et pour $i$ quelconque, % L7
on définit % L8
\[
M\in\mathrm{Ob}\,\mathcal M^i(X)\Longleftrightarrow M(-j)\in\mathrm{Ob}\,\mathcal M^{i+2j}(X) % L9
\]
où on prend $j$ assez grand pour que % L10
$M(-j)\in\mathrm{Ob}\,M^+(X)$. On notera que % L11
\[
\mathcal M^{+i}(X)=\mathcal M^i(X)\cap\mathcal M^+(X) % L12
\]
% slanted
\margin{indices en bas} % L12a
si $i\geq0$, et pour tt $i$ si on définit % L13
$\mathcal M^{\unsure{+i|-i}}(X)=\{0\}$\deleted{\unsure{$\mathcal M^+(X)$}} si $i<0$ \deleted{(\unsure{ce qui revient au|ce --- pour les} % L14
même si on veut que $\mathcal M^{+i}(X)$ \inserted{$\subset\mathcal M^+(X)$ et} \unsure{existe aussi ici|consiste aussi ici}}] % L15

On aura % L16
\[
\begin{aligned}
\otimes\deleted{\mathbb Z}\mathbb Q(-1)&:\quad\mathcal M^i(X)\overset{\text{équiv}}{\simeq}\mathcal M^{i+2}(X)\\ % L17
\otimes\deleted{\mathbb Z}\mathbb Q(j)&:\quad\mathcal M^i(X)\overset{\text{équiv}}{\simeq}\mathcal M^{\unsure{i-2j|i+2j}}(X) % L18
\end{aligned}
\]

% archive sheet number: 32
\scanpage{34}
mais on fera attention que l'inclusion % L1
\[
\begin{aligned}
\deleted{\mathbb Z}\mathbb Q(-1)\otimes\mathcal M^{+i}(X)&\hookrightarrow\mathcal M^{+(i+2)}(X)\qquad\text{plus gén.}\\ % L2
\deleted{\mathbb Z}\mathbb Q(-j)\otimes\mathcal M^{+i}(X)&\unsure{\lhook\joinrel-|\hookrightarrow}M^{+(i+2j)}(X) % L3
\end{aligned}
\]
est \uline{stricte} en général, i.e. n'est pas % L4
une équivalence de catégories. % L5

Noter que $\mathcal M^{+i}(X)$ est \uline{épaisse} dans $\mathcal M^{+j}(X)$ % L6
% slanted
\margin{indices en bas} % L6a
\deleted{[elle est définie, dans la catégorie}, et de % L7
$=$ $\mathcal M^i(X)$ épaisse dans $\mathcal M^j(X)$ $(i\leq j)$. % L8

Les quotients $G_i^+(X)=\mathrm{Gr}^{\unsure{F|\mathscr F}}_i(\mathcal M^+(X)\deleted{\unsure{)_{i-1}}})\inserted{\simeq\mathcal M_i^+(X)/\mathcal M_{i-1}^+(X)}$ et $G_i(X)=\mathrm{Gr}^{F}_i(\mathcal M(X))$ % L9
$\simeq\mathcal M_i(X)/\mathcal M_{i-1}(X)$ sont fort intéressants. % L10

Noter que les \deleted{\unsure{M}}$G_i(X)$ sont donc \deleted{\unsure{isomorphes}} \inserted{équivalentes} % L11
par twisting $G_i(X)\overset{\mathbb Z(-j)}{\simeq}G_{i+2j}(X)$, en % L12
particulier tous \deleted{\unsure{isomorphes}} \inserted{équiv. can.} à $G_0(X)$. % L13
D'ailleurs $\unsure{G_i^+|G\mathrm r_i^+|\mathcal G_i^+}(X)\underset{\sim}{\hookrightarrow}G_i(X)\overset{\sim}{\simeq}G_0(X)$ % L14
\inserted{équiv. à une} sous catégorie pleine et \uline{épaisse} de $G_i(X)\overset{\sim}{\simeq} % L15
G_0(X)$. % L16

% archive sheet number: 33
\scanpage{35}
b) filtration par ``\deleted{la}\unsure{le|la} \deleted{dimension}'' ``type dimensionnel''. % L1

On pose pour $i\geq 0$ % L2
\[\mathcal D_i(\mathcal M^+(X))=\deleted{\unsure{\mathcal M|\mathcal M^i}}\,\text{sous-catégorie pleine}\] % L3
de $\mathcal M^+(X)$ formée des objets % L4
qui se \inserted{loc. étale} dévissent [(du moins % L5
fibre par fibre)] en objets % L6
\margin{\deleted{\unsure{$M$|$\mathrm{Ob}$}}} % L6a
de la forme % L7
\deleted{$\unsure{M}$} $M(-j)$, avec $M\in\mathrm{Ob}\,\mathcal M^{+i}(X)$, % L8
\inserted{$j\in\mathbb Z,\ j\geq 0$} % L8a
i.e. on prend (du moins pour $X=\mathrm{Spec}\ k$ d'un % L9
corps) la filtration minima qui % L10
majore celle par le poids, \deleted{\unsure{et}} qui soit % L11
stable par $\otimes\deleted{\mathbb Z}\mathbb Q(-1)$, et \uline{épaisse}. % L12

Les $\mathcal D_i(\mathcal M^+(X))$ sont stables par l'image inverse, % L13
et pour l'image directe on a % L14
\[\begin{aligned}\mathbb R^j f_!(\mathcal D_i(\mathcal M^+(X)))&\subset\mathcal D_{i+j}(\mathcal M^+(Y))\\ % L15
\mathbb R^j f_*(\mathcal D_i(\mathcal M^+(X)))&\subset\mathcal D_{i+j}(\mathcal M^+(Y))\ \text{\textcircled{?}}\end{aligned}\] % L16
% margin note written sideways
\margin{\deleted{\unsure{le cas critique|les cas critiques} \unsure{est celui d'une} immersion ouverte}} % L16a
% slanted
\margin{cela distingue formellement les filtrations $\mathcal D_i(M)$ et $M_i$ !} % L16b
\inserted{[On a seulement $\mathbb R^jf_*(\mathcal M^{+(i)}(X))\subset\mathcal M^{+(i+2j)}(Y)$ !?]} % L17
% slanted
\margin{\unsure{tout reformuler|tout se formule|toute formule} en termes de la dim. de $X$ !} % L18

Ici \deleted{\unsure{encore}} on a % L19
\[M\in\mathrm{Ob}\,\mathcal D_i(\mathcal M^+(X))\Longleftrightarrow M(-j)\in\deleted{\unsure{\mathrm{Ob}}}\mathcal D_{i\deleted{\unsure{+j}}}(\mathcal M^+(X))\] % L20
\margin{pas $\llap{\uuline{\phantom{pas}}\kern.25em}\mathcal D_{i+j}$ !!} % L20a
et la filtration \inserted{de $\mathcal M^+(X)$} par les $\mathcal D_i(\mathcal M^+(X))$ s'étend % L21
en une filtration de $\mathcal M(X)$ par des $\mathcal D^i(\mathcal M(X))$. % L22

% archive sheet number: 34
